\section{Method}
\label{sec:method}

\subsection{Overview}

MGR-SID keeps MiniOneRec's semantic \rqvae{} tokenizer and augments it with a train-only graph bank. The core intuition is simple: different SID levels solve different structural discrimination problems, so they should preserve different collaborative graph structures. In the current \textbf{v1} implementation, we instantiate this idea with a fixed but non-uniform level-to-view assignment:
\begin{itemize}[leftmargin=1.4em]
  \item level 1 is regularized by a coarse collaborative graph;
  \item level 2 is regularized by a middle-resolution spectral graph;
  \item level 3 is regularized by a local transition graph.
\end{itemize}

This is intentionally simpler than the full long-term vision of learned level-wise allocation. The purpose of v1 is to test whether hierarchy-aware structural supervision is already useful before adding another learned allocation module.

\subsection{Semantic backbone}

Given an item embedding $\mathbf{x}_i$, the encoder produces a latent representation that is quantized by three residual quantizers. Let $\mathbf{q}_i^{(l)}$ be the quantized vector produced by the $l$-th quantizer, and let
\begin{equation}
\hx_i^{(l)} = \sum_{r=1}^{l} \mathbf{q}_i^{(r)}
\end{equation}
be the cumulative quantized representation after level $l$. The decoder reconstructs $\mathbf{x}_i$ from $\hx_i^{(3)}$. This is exactly the same semantic tokenizer backbone used by MiniOneRec, and we keep its reconstruction loss and residual quantization loss unchanged.

\subsection{Graph bank construction}

All graph views are constructed from training interactions only. For each interaction sample with history $H_t = [i_{t-k}, \dots, i_{t-1}]$ and target item $y_t$, we build three graph carriers.

\paragraph{Coarse collaborative graph.}
We build an undirected item-item graph from the unique items appearing in the recent history window and the target. For items $u$ and $v$ that co-occur inside the sequence, we add a symmetric edge weighted by inverse distance. This graph captures broad compatibility and stable collaborative organization. We then purify it by support pruning, popularity debiasing, and row normalization:
\begin{equation}
\Acoarse = \mathrm{RowNorm}\!\left( \mathrm{Debias}\!\left( \mathrm{Prune}(\Gcoarse) \right)\right).
\end{equation}
In the current implementation, pruning removes edges with weight below $2.0$ and popularity debiasing uses exponent $\alpha=0.5$.

\paragraph{Local transition graph.}
We build a directed graph from recent history items to the target item. Each edge receives inverse-rank weight, so more recent items contribute more. The purified local graph applies support pruning, target-popularity correction, and row normalization:
\begin{equation}
\Alocal = \mathrm{RowNorm}\!\left( \mathrm{TargetDebias}\!\left( \mathrm{Prune}(\Glocal) \right)\right).
\end{equation}
This view captures short-range next-item preference but is naturally sparse and noisy.

\paragraph{Middle-resolution graph.}
The middle view is the key design choice. Early heuristic probes using diffusion residuals and simple band-pass proxies already suggested that mid-scale structure matters, but they were weaker than spectral graph constructions. The current v1 therefore selects a FaGSP-inspired spectral middle view as the default:
\begin{equation}
\Amid = \mathrm{BandReconstruct}\!\left( \mathrm{SymNorm}(\Acoarse)\right).
\end{equation}
Operationally, we compute a band-limited spectral reconstruction of the purified coarse graph by retaining eigen-components in a middle spectral range. The implementation uses a rank of $48$ and an eigen-ratio range of $[0.25, 0.65]$. We also explored a GSPRec-style temporal spectral view as an alternative, but use the FaGSP-style view as the main $G_{\mathrm{mid}}$ because it performed best in our diagnostic probes.

\paragraph{Semantic-anchor purification.}
Inspired by PRISM, we additionally tested a semantic-anchor purification stage in which a purified graph is intersected with a top-$k$ semantic $k$NN graph extracted from the text embedding space. This module improves some graph views, but our probes indicate that the main gain comes from the quality of the middle-resolution graph itself rather than from anchor-only denoising. In the current training-time v1, the final chosen middle view is the spectral graph built from the purified coarse graph.

\subsection{Hierarchy-aware graph regularization}

The semantic backbone already optimizes
\begin{equation}
\Lrec + \Lrq,
\end{equation}
where $\Lrec$ is the reconstruction loss and $\Lrq$ is the residual quantization loss. MGR-SID v1 adds graph regularization on the cumulative representations $\hx_i^{(l)}$.

For a mini-batch $\mathcal{B}$ and a row-normalized subgraph adjacency $\mathbf{A}_{\mathcal{B}}^{(l)}$ extracted from the selected graph view, we define
\begin{equation}
\Lgraph^{(l)} =
\frac{1}{|\mathcal{B}|}
\sum_{i \in \mathcal{B}}
\left\|
\hx_i^{(l)} -
\sum_{j \in \mathcal{B}} \mathbf{A}^{(l)}_{\mathcal{B},ij} \hx_j^{(l)}
\right\|_2^2 .
\label{eq:graph_smooth}
\end{equation}
This loss favors level-$l$ representations that are locally smooth on the corresponding graph view. Intuitively, items that are connected in the relevant graph should remain nearby at the level that is responsible for preserving that structure.

The hierarchy-aware v1 objective is then
\begin{equation}
\Ltotal =
\Lrec + \Lrq
 + \lambda_1 \Lgraph^{(1)}
 + \lambda_2 \Lgraph^{(2)}
 + \lambda_3 \Lgraph^{(3)},
\end{equation}
where
\begin{equation}
\mathbf{A}^{(1)}_{\mathcal{B}} \leftarrow \Acoarse,\quad
\mathbf{A}^{(2)}_{\mathcal{B}} \leftarrow \Amid,\quad
\mathbf{A}^{(3)}_{\mathcal{B}} \leftarrow \Alocal .
\end{equation}
In the current implementation, $(\lambda_1,\lambda_2,\lambda_3) = (0.05, 0.15, 0.05)$.

\subsection{Controls}

To isolate the effect of hierarchy-awareness, we compare MGR-SID v1 with two controls:
\begin{itemize}[leftmargin=1.4em]
  \item \textbf{semantic baseline}: the upstream-aligned MiniOneRec \rqvae{} tokenizer with no graph regularization;
  \item \textbf{uniform regularization}: a control that averages the coarse, middle, and local graphs and applies one graph smoothness loss to the final cumulative representation only.
\end{itemize}
The uniform control keeps the same graph ingredients but removes the hierarchy-aware assignment. It is therefore the most direct test of whether the benefit comes from non-uniform level-specific supervision rather than from adding graph information in any form.

\subsection{Current scope}

MGR-SID v1 should be read as a tokenizer-stage instantiation, not the final method ceiling. It does \emph{not} yet learn the graph allocation itself, and it does \emph{not} modify downstream SFT or RL. The purpose of the current draft is narrower: to establish whether graph-structured collaborative information helps the semantic tokenizer \emph{when injected as level-specific structural supervision}.
